%%  Moved from the article's Threats to validity (80_limits.tex) in the IJF
%%  retarget, PLAN-IJF.md Task 2.1; the article keeps each paragraph's
%%  verdict and points here for the argument.

\paragraph{What the band's coverage does and does not establish --- the three qualifications}
The measurement switches the resampling off, so it
bounds the critical value's contribution and not the bootstrap's: the
pointwise interval underneath is governed by $K/n$ and undercovers where
this corpus lives (Section~\ref{sec:simsensitivity}), and the $(n, K)$
calibration that would widen it is reported as a sensitivity and not
applied, because it is fitted pointwise, applied simultaneously and
measurably misspecified. The level is bought with width:
\nResolvedLostToEmp\ cells the multiplier would have resolved are unresolved
here, and \nUnresolvedPairs\ pairs resolve nothing. And the decision-curve
families carry cells that are zero by arithmetic, which a declared rule
excludes from the family before the quantile is taken
(Section~\ref{sec:dcabands}). The decision-curve region of that section is
the exception to the first sentence of this paragraph's verdict: its family's
empirical quantile covers \covEmpHeavyCaseAtDesign\ on a matched family that
reproduces the corpus's estimator ratio at \ratioSynthDca\ against
\ratioCorpusDca, so that label rests on a critical value about a point short
of its level.

\paragraph{Internal --- two mechanisms this evidence retires}
A multinomial resample holds about $1 - e^{-1}$ of a high-cardinality
register's levels, which is the natural explanation of a displaced bootstrap
distribution. Under weights every row has positive weight in every draw and
the level share is \levelShareWeighted\ against \levelShareMult\ --- the
mechanism is gone by construction --- and the displacement's median moves
only from \dispMedianMult\ to \dispMedianWeighted\ (Section~\ref{sec:boot}).
What the change does remove is a defect the mechanism genuinely caused:
cells a draw drops entirely, which get no band at all and were counted as
\emph{unresolved} beside cells that were banded and straddled zero ---
\nUnbandableSameDesign\ under the multinomial scheme on this design,
\nUnbandableCells\ under weights. And the draw count is not what governs the
multiplier's coverage: on families matched to this design it is
\covMultBootOneFifty\ at \nBootOneFifty\ draws, \covMultBootFourHundred\ at
\nBootFourHundred\ and \covMultBootOneThousand\ at \nBootOneThousand\ ---
it plateaus --- while the empirical quantile's coverage is what rises with
$B$, from \covEmpBootThirtyThree\ at \nBootThirtyThree\ draws to
\covEmpBootFourHundred\ at \nBootFourHundred. A mechanism can be removed
completely and leave its supposed consequence standing, which is a caution
about every explanation of a resampling artefact offered without the
counterfactual, including the two here.

\paragraph{Internal --- further assumptions, named, and one of them
measured}
The split point is fixed within a draw, so split-position uncertainty is
outside every interval here. A second sort is undeclared and is not repaired:
the register-quality mechanisms cut a cumulative-count curve inside a group
of tied identities, so \emph{which} identities a degradation removes rests on
an order nothing pins; the affected quantity is the quality axis and not the
increment at the reference cell, and repairing it moves every number on that
axis, which is why it is recorded here and why
Supplement~\ref{app:standard}'s tie-break item ranges over every sort a
procedure cuts on. The moving-block construction needs stationarity and
mixing that an event log is not guaranteed to have, and the same assumption
makes Definition~\ref{def:estimand}'s population limit well defined under a
temporal split. Stationarity is measured and does not hold everywhere:
cutting each pair's test half into \nDriftBlocks\ consecutive blocks and
reading the spread of the increment against a null that re-partitions the
same rows at random, \nDrifts\ of \nDriftPairs\ pairs sit in the upper tail
of that null at $\alpha = 0.05$, \nDriftSignVaries\ carry both signs between
blocks of one test half, and only \nDriftTrend\ shows a monotone trend, so
this is heterogeneity and not decay (Supplement~\ref{app:drift}). On those
pairs the interval and the band understate what a reader would see on a
different stretch of the same log, by an amount this design cannot sign;
time is not a declared axis of $\mathcal{S}$. Mixing is not tested. The
band's family-wise coverage is asymptotic, the block length is a convention,
$\lceil n^{1/3}\rceil$, which moves coverage by up to
\coverageBlockSpreadMax\ within a world, and the $(n, K)$ calibration of
Supplement~\ref{app:calibration} corrects the studentised statistic's
\emph{scale} and not the shape of the maximum's distribution and is
measurably misspecified ($\log n$ enters at $t = \planeNdependenceT$). A
reader who rejects any of these should rely on the objects that need no
band: the point estimates, the decomposition, the disagreement rates and the
regret comparison.

